\begin{abstract}
Humans write code in a fundamentally interactive manner and rely on constant execution feedback to correct errors, resolve ambiguities, and decompose tasks. While LLMs have recently exhibited promising coding capabilities, current coding benchmarks mostly consider a static instruction-to-code sequence transduction process, which has the potential for error propagation and a disconnect between the generated code and its final execution environment. To address this gap, we introduce \benchmark{}, a lightweight, flexible, and easy-to-use framework of interactive coding as a standard reinforcement learning (RL) environment, with code as actions and execution feedback as observations. Our framework is language and platform agnostic, uses self-contained Docker environments to provide safe and reproducible execution, and is compatible out-of-the-box with traditional seq2seq coding methods, while enabling the development of new methods for interactive code generation.
% \benchmark{}'s execution-oriented evaluation offers a multitude of novel feedback signals that can be used to develop interactive coding agents that leverage the extra contextual information.
% \benchmark{} defines a general formulation for \textit{interactive coding tasks} and provides a framework that makes it easy to construct such a task setting from existing, static instruction-to-code datasets.
We use \benchmark{} to create three interactive code environments with Bash, SQL, and Python as action spaces, leveraging data from the static NL2Bash~\cite{lin-etal-2018-nl2bash}, Spider~\cite{yu-etal-2018-spider}, and MBPP~\cite{austin2021program} datasets. We demonstrate \benchmark{}'s viability as a testbed by evaluating multiple state-of-the-art LLMs configured with different prompting strategies such as ReAct~\cite{yao2023react} and Plan \& Solve~\cite{wang2023planandsolve}. Our results showcase the benefits of interactive code generation and demonstrate that \benchmark{} can serve as a challenging benchmark for advancing code understanding and generation capabilities. \benchmark{} is designed to be easily extensible and can even be used to create new tasks such as Capture the Flag, a popular coding puzzle that is inherently multi-step and involves multiple programming languages.
\footnote{Code and data available at \url{https://intercode-benchmark.github.io/}}
% % Finally, we develop a new Capture the Flag task with \benchmark{} that showcases \benchmark{} value towards constructing new tasks, datasets, and evaluation metrics.
% \ap{
% that showcases the potential of \benchmark{} in enabling the development of innovative tasks and metrics.}
\end{abstract}